% chapters/journal/04-method.tex
% 第4章 方法设计

\section{方法设计}

\subsection{交易特征提取}

对每笔候选交易 $tx_i$，提取 $d$ 维原始特征向量 $\mathbf{x}_i$ 后，通过两层全连接网络编码为稠密表示：

\begin{equation}
  \mathbf{h}_i = \text{ReLU}(\mathbf{W}_2 \cdot \text{ReLU}(\mathbf{W}_1 \mathbf{x}_i + \mathbf{b}_1) + \mathbf{b}_2)
\end{equation}

其中 $\mathbf{h}_i \in \mathbb{R}^{d_h}$ 为交易 $tx_i$ 的编码表示，$d_h$ 为隐藏层维度。交易特征编码网络的参数在策略网络训练过程中联合优化。

原始特征向量 $\mathbf{x}_i$ 由三组特征拼接而成。交易基础特征包括归一化手续费、Gas 消耗估计、到达时间戳和交易类型的独热编码。历史统计特征包括发送地址的历史交易成功率和历史确认延迟。风险特征包括基于历史异常行为统计生成的风险评分，该评分综合考虑发送地址参与异常交易（如失败率过高、高频提交）的历史频率。

\subsection{策略网络结构}

本文采用近端策略优化（Proximal Policy Optimization，PPO）\cite{schulman2017ppo}算法训练排序策略。选择 PPO 而非 DQN\cite{mnih2015dqn} 的原因在于：本文的动作空间随候选交易集合的变化而动态改变，传统 DQN 的固定输出维度难以适应这一特性；PPO 基于策略梯度的框架通过对每笔候选交易独立计算选择得分，天然支持变长动作空间。策略网络采用 Actor-Critic 架构，Actor 和 Critic 共享交易特征编码层，在此基础上分别接各自的决策头。

Actor 网络接收状态 $s_t$，对每笔候选交易计算选择得分：

\begin{equation}
  u_i = \mathbf{w}_a^\top \text{ReLU}(\mathbf{W}_s [\mathbf{h}_i; \mathbf{b}_t] + \mathbf{b}_s)
\end{equation}

其中 $[\mathbf{h}_i; \mathbf{b}_t]$ 为交易表示与区块状态的拼接向量。选择得分经 softmax 归一化后得到动作概率分布 $\pi_\theta(a|s_t)$。

Critic 网络估计当前状态的价值函数 $V_\psi(s_t)$，用于计算优势函数以降低策略梯度的方差。Critic 网络对所有候选交易表示取均值池化后与区块状态拼接，经两层全连接网络输出标量价值估计：

\begin{equation}
  V_\psi(s_t) = \mathbf{w}_v^\top \text{ReLU}(\mathbf{W}_v [\bar{\mathbf{h}}_t; \mathbf{b}_t] + \mathbf{b}_v) + b_v
\end{equation}

其中 $\bar{\mathbf{h}}_t = \frac{1}{|\mathcal{T}_t|}\sum_{tx_i \in \mathcal{T}_t} \mathbf{h}_i$ 为候选交易表示的均值池化。与 Actor 网络需要对每笔交易分别计算得分不同，Critic 网络通过均值池化将变长的候选集合压缩为固定维度的向量，从而输出对当前状态的整体价值估计。

\subsection{合法性掩码机制}

由于候选交易集合的大小在不同时间步动态变化，且部分交易受容量约束和依赖约束限制不可被选取，本文引入合法性掩码（action masking）机制。在计算动作概率分布时，将不属于合法动作集合 $\mathcal{A}_t^{\text{valid}}$ 的交易得分设为负无穷，再进行 softmax 归一化：

\begin{equation}
  \pi_\theta^{\text{mask}}(a_i|s_t) = \frac{\exp(u_i) \cdot \mathbb{I}(a_i \in \mathcal{A}_t^{\text{valid}})}{\sum_{a_j \in \mathcal{A}_t^{\text{valid}}} \exp(u_j)}
\end{equation}

该机制确保策略网络仅在合法动作上分配概率，避免生成违反约束的无效排序。

为支持结构消融实验，本文在实现中进一步提供两类可控开关：其一，关闭动作掩码（No-ActionMask）以检验合法性先验对策略学习的贡献；其二，关闭 \texttt{STOP} 动作（No-STOP）以检验主动终止能力对风险控制和打包效率的影响。完整模型保留动作掩码与 \texttt{STOP}，作为默认配置。

\subsection{训练算法}

训练过程采用 PPO 的截断替代目标函数。在每个训练回合中，智能体与仿真环境交互完成一次完整的区块构建过程，收集轨迹 $\tau = \{(s_t, a_t, r_t, s_{t+1})\}_{t=0}^{T}$。即时奖励按第3章定义计算，即手续费收益、增量等待时间均衡奖励、位置敏感风险惩罚与提前终止惩罚的加权组合。使用广义优势估计（Generalized Advantage Estimation，GAE）计算优势函数：

\begin{equation}
  \hat{A}_t = \sum_{l=0}^{T-t} (\gamma \lambda_{\text{GAE}})^l \delta_{t+l}
\end{equation}

其中 $\delta_t = r_t + \gamma V_\psi(s_{t+1}) - V_\psi(s_t)$ 为时序差分残差，$\gamma$ 为折扣因子，$\lambda_{\text{GAE}}$ 为 GAE 参数。

策略网络的优化目标为 PPO 截断目标：

\begin{equation}
  L^{\text{CLIP}}(\theta) = \mathbb{E}_t \left[ \min\left( \rho_t \hat{A}_t, \; \text{clip}(\rho_t, 1-\epsilon, 1+\epsilon) \hat{A}_t \right) \right]
\end{equation}

其中 $\rho_t = \pi_\theta(a_t|s_t) / \pi_{\theta_{\text{old}}}(a_t|s_t)$ 为重要性采样比率，$\epsilon$ 为截断参数。Critic 网络通过最小化价值函数的均方误差进行更新：

\begin{equation}
  L^{\text{value}}(\psi) = \mathbb{E}_t \left[ (V_\psi(s_t) - \hat{R}_t)^2 \right]
\end{equation}

其中 $\hat{R}_t = \sum_{l=0}^{T-t} \gamma^l r_{t+l}$ 为从时间步 $t$ 开始的折扣累积奖励。训练过程中的主要超参数配置如表~\ref{tab:hyperparams} 所示。

\begin{table}[htbp]
\centering
\caption{训练超参数配置}
\label{tab:hyperparams}
\begin{tabular}{llc}
\toprule
超参数 & 符号 & 取值 \\
\midrule
Actor 学习率 & $\alpha_\theta$ & $3 \times 10^{-4}$ \\
Critic 学习率 & $\alpha_\psi$ & $1 \times 10^{-3}$ \\
折扣因子 & $\gamma$ & 0.99 \\
GAE 参数 & $\lambda_{\text{GAE}}$ & 0.95 \\
PPO 截断参数 & $\epsilon$ & 0.2 \\
每回合更新轮数 & $E$ & 4 \\
隐藏层维度 & $d_h$ & 128 \\
训练回合数 & $M$ & 3{,}000 \\
\bottomrule
\end{tabular}
\end{table}

训练算法的完整流程如算法~\ref{alg:training} 所示。

\begin{algorithm}[H]
\caption{基于 PPO 的交易排序策略训练}
\label{alg:training}
\begin{algorithmic}[1]
\REQUIRE 仿真环境 $\text{Env}$, 学习率 $\alpha_\theta, \alpha_\psi$, 截断参数 $\epsilon$
\ENSURE 训练后的策略参数 $\theta^*$
\STATE 初始化 Actor 参数 $\theta$, Critic 参数 $\psi$
\FOR{episode $= 1, 2, \ldots, M$}
    \STATE 从环境采样候选交易集合 $\mathcal{T}$，初始化区块状态
    \STATE 收集完整轨迹 $\tau = \{(s_t, a_t, r_t)\}_{t=0}^{T}$
    \STATE 计算 GAE 优势估计 $\hat{A}_t$
    \FOR{epoch $= 1, 2, \ldots, E$}
        \STATE 更新 $\theta \leftarrow \theta + \alpha_\theta \nabla_\theta L^{\text{CLIP}}(\theta)$
        \STATE 更新 $\psi \leftarrow \psi - \alpha_\psi \nabla_\psi L^{\text{value}}(\psi)$
    \ENDFOR
\ENDFOR
\RETURN $\theta^*$
\end{algorithmic}
\end{algorithm}

推理阶段，策略网络对候选交易计算选择得分并施加合法性掩码后，取概率最高的交易作为下一步选择（贪心解码），依次构建完整的区块执行序列。

\subsection{推理过程}

训练完成后，策略网络以贪心方式执行推理。给定一组候选交易集合 $\mathcal{T}$ 和初始空区块状态，推理过程按以下步骤逐笔构建区块执行序列：

（1）计算当前时间步所有候选交易的选择得分，并施加合法性掩码过滤不可选交易。

（2）从合法动作集合中选取得分最高的交易 $a_t^* = \arg\max_{a \in \mathcal{A}_t^{\text{valid}}} \pi_\theta^{\text{mask}}(a|s_t)$ 加入区块执行序列。若 \texttt{STOP} 动作得分最高，则终止当前区块构建。

（3）更新区块状态，重建合法动作集合，进入下一时间步。

（4）重复上述过程直至满足终止条件：合法候选交易为空、剩余 Gas 不足、达到最大选择步数、或智能体选择 \texttt{STOP}。

推理结束后输出完整的交易执行序列 $\sigma = (tx_{\pi(1)}, \ldots, tx_{\pi(K)})$。与训练阶段的随机采样不同，推理阶段采用确定性贪心选择，以保证评估结果的可复现性。

\subsection{消融设计对应的结构说明}

为验证模型关键组件的必要性，本文设计四个结构变体用于消融实验：

（1）\textbf{No-SeqSummary}：移除区块状态向量中的已选交易序列摘要 $\mathbf{h}_t^{\text{seq}}$，Actor 和 Critic 网络仅接收交易特征与标量区块状态（剩余 Gas、已选数量、累计收益）作为输入。该变体用于检验序列上下文信息对逐笔决策质量的贡献。

（2）\textbf{No-ActionMask}：关闭合法性掩码机制，策略网络在全部候选交易（含不合法交易）上计算概率分布。当智能体选择不合法交易时，环境不执行该动作并施加负奖励惩罚，但不会因此进入无限循环。该变体用于评估合法动作先验对策略学习效率和最终性能的影响。

（3）\textbf{No-STOP}：移除 \texttt{STOP} 动作，智能体只能通过外部终止条件（候选池耗尽、Gas 不足、达到步数上限）结束区块构建。该变体用于检验主动终止能力对风险控制和打包效率的贡献。

（4）\textbf{Full Model}：保留所有组件（序列摘要、动作掩码、\texttt{STOP} 动作）的完整模型，作为消融实验的对照基准。
